\documentclass[11pt,a4paper]{ctexart}

% === 页面设置 ===
\usepackage[margin=2.5cm]{geometry}
\usepackage{hyperref}
\usepackage{graphicx}
\usepackage{amsmath,amssymb}
\usepackage{booktabs}
\usepackage{xcolor}
\usepackage{enumitem}
\usepackage{caption}
\usepackage{fancyhdr}

% === 代码高亮 ===
\usepackage{listings}
\definecolor{codebg}{RGB}{245,245,245}
\definecolor{codeframe}{RGB}{200,200,200}
\definecolor{commentgreen}{RGB}{0,128,0}
\definecolor{keywordblue}{RGB}{0,0,180}
\definecolor{stringred}{RGB}{180,0,0}

\lstset{
  backgroundcolor=\color{codebg},
  frame=single,
  framerule=0.5pt,
  rulecolor=\color{codeframe},
  basicstyle=\ttfamily\small,
  breaklines=true,
  breakatwhitespace=false,
  columns=flexible,
  keepspaces=true,
  tabsize=4,
  showstringspaces=false,
  commentstyle=\color{commentgreen}\itshape,
  keywordstyle=\color{keywordblue}\bfseries,
  stringstyle=\color{stringred},
  numbers=left,
  numberstyle=\tiny\color{gray},
  numbersep=8pt,
  xleftmargin=2em,
  aboveskip=1em,
  belowskip=1em,
  language=Python,
  morekeywords={dtype,float32,uint8,int64,arange,meshgrid,zeros_like,where,exp,sqrt,cos,sin,pi,abs,angle,mean,std,clip,power,maximum,fftshift,ifftshift,fft2,ifft2,arctan2,copy,sum},
  escapeinside={(*@}{@*)},
}

% === 页眉页脚 ===
\pagestyle{fancy}
\fancyhf{}
\fancyhead[L]{\small\textbf{代码导读 — RTAD MRAF/GS 管线}}
\fancyhead[R]{\small\thepage}
\renewcommand{\headrulewidth}{0.5pt}

% === 标题 ===
\title{\textbf{代码导读\\[0.3em]\large 从初始相位到 SLM 输出的核心算法逐段解释}}
\author{RTAD MRAF/GS Pipeline}
\date{\today}

\begin{document}

\maketitle
\thispagestyle{fancy}

\vspace{1em}

\begin{abstract}
\noindent 读完这份文档，你就能理解从初始相位到 SLM 输出的每一行关键代码在做什么。所有代码片段均从本目录的 \texttt{.py} 文件中直接提取。
\end{abstract}

\tableofcontents
\newpage

% ======================================================================
\section{管线数据流}

整个系统分为三个阶段：

\begin{center}
\begin{tabular}{ccccc}
\texttt{make\_phase0.py} & & \texttt{run\_rtad\_mraf\_gs\_case.py} & & \texttt{convert\_to\_slm.py} \\
(Stage 1) & & (Stage 2) & & (Stage 3) \\[4pt]
beam=3.8mm & & \texttt{phase0.mat} & & $2048\times2048$ 相位 \\
$\downarrow$ & & $\downarrow$ & & $\downarrow$ \\
Romero-Dickey & $\longrightarrow$ & WGS 迭代 200轮 & $\longrightarrow$ & 复振幅插值 $\rightarrow$ SLM BMP \\
解析相位 & & (MRAF投影+权重更新) & & $1024\times1024$ @ 17$\mu$m \\
\end{tabular}
\end{center}

\vspace{1em}
\textbf{三个 Stage 的关系：}
\begin{itemize}[nosep]
  \item \textbf{Stage 1} 生成初始相位猜测（Romero-Dickey 解析法），输出 \texttt{phase0.mat}。
  \item \textbf{Stage 2} 用 MRAF + WGS 迭代 200 轮优化相位，记录收敛曲线和指标。
  \item \textbf{Stage 3} 把 $2048\times2048$ 的优化相位插值缩放到 $1024\times1024$，叠加闪耀光栅，输出 SLM 所需的 BMP 文件。
\end{itemize}

% ======================================================================
\section{目标生成 — $330\times120\mu\text{m}$ 矩形平顶}

\textbf{文件}: \texttt{src/rtad\_target.py}

\subsection{升余弦边缘}

矩形不是硬边界，而是通过升余弦 (raised cosine) 从 flat 区平滑过渡到零:

\begin{lstlisting}
# src/rtad_target.py, line 69-78
def raised_cosine_edge(u, u0, u1):
    """一维升余弦下降沿: u<=u0->1, u>=u1->0, 中间余弦过渡."""
    C = np.zeros_like(u, dtype=np.float32)
    C[u <= u0] = 1.0                    # flat 区 = 1
    idx = (u > u0) & (u < u1)
    t = (u[idx] - u0) / (u1 - u0)      # 归一化到 [0,1]
    C[idx] = 0.5 * (1.0 + np.cos(np.pi * t))  # 余弦曲线从 1 降到 0
    return C
\end{lstlisting}

\textbf{物理含义}: 这是一维 profile。对于 $330\times120\mu\text{m}$ 的目标:
\begin{itemize}[nosep]
  \item $a_{50} = 165\mu\text{m}$ (X 方向半宽, 50\% 强度处)
  \item $a_0 = a_{50} - \delta_x = 150\mu\text{m}$ (flat 区边界, 强度=1)
  \item $a_1 = a_{50} + \delta_x = 180\mu\text{m}$ (过渡区外边界, 强度=0)
  \item $\delta_x = 15\mu\text{m}$ 控制边缘过渡的陡峭程度
\end{itemize}

二维模板 $I(x,y) = C_x(|x|) \times C_y(|y|)$:

\begin{lstlisting}
# src/rtad_target.py, line 151-154
Ix = raised_cosine_edge(absX, a0, a1)
Iy = raised_cosine_edge(absY, b0, b1)
I_full = np.clip(Ix * Iy, 0.0, 1.0)   # 二维强度模板 (可分离)
A_full = np.sqrt(I_full)               # 振幅 = sqrt(强度)
\end{lstlisting}

\subsection{四层 Mask 区域}

MRAF 算法的核心是把焦面分成四个区，每个区有不同约束规则:

\begin{lstlisting}
# src/rtad_target.py, line 161-171
# 从内到外四个区域:
mask_flat   = (absX <= a0) & (absY <= b0)          # 平顶核心 (150x52 um)
mask_signal = I_full >= release_level               # 约束区 (I >= e^{-2})
mask_signal = mask_signal | mask_flat               # 确保 flat 一定在 signal 内
mask_free   = mask_guard_window & ~mask_signal      # 自由区 (可随意演化)
mask_bg_far = ~(mask_signal | mask_free)            # 远背景 (轻微衰减)
\end{lstlisting}

\textbf{区域示意图} (以 X 方向为例):

\begin{verbatim}
强度
1.0 |  xxxxxxxxxxxx
    |  x  flat   x\
0.5 |  x  (a0=150) x\___ 升余弦过渡 ___
    |  x           x        \
e^{-2} |  xxxxxxxxxxxxxxxxxxxxxx\_________
    |  |<-- signal -->|              \
0   |  |<--          guard_window         -->|
    |  |<--    free    -->|<--  bg_far  -->|

 <-- 区域边界: a0=150, a50=165, a1=180, a2=200um -->
\end{verbatim}

% ======================================================================
\section{初始相位 — Romero-Dickey 解析法}

\textbf{文件}: \texttt{make\_phase0.py}

DOE 初始相位不能用随机值——那样迭代根本收不敛。Romero-Dickey 方法从高斯光束和矩形平顶的能量守恒关系出发，解析求出可分离的初始相位:

\begin{lstlisting}
# make_phase0.py, line 44-51
def romero_dickey_phase_1d(x_m, ri_m, Ro_m, lambda_m, f_m):
    """Romero-Dickey 一维解析相位."""
    xi = x_m / ri_m                           # 归一化坐标 (ri=1/e 振幅半径)
    # 无量纲相位函数:
    phi_dimless = (xi * np.sqrt(np.pi) / 2.0 * erf(xi)
                   + 0.5 * np.exp(-xi**2) - 0.5)
    # beta: 无量纲参数, 表示几何光学近似程度
    beta = 2.0 * np.pi * ri_m * Ro_m / (lambda_m * f_m)
    # beta x phi_dimless -> 真实相位 (rad)
    return beta * phi_dimless
\end{lstlisting}

\textbf{$\beta$ 的物理含义}:
\[
\beta = 2\pi \times \frac{(\text{光束半径}) \times (\text{目标半宽})}{\lambda \times (\text{焦距})}
\]
\begin{itemize}[nosep]
  \item $\beta > 10 \rightarrow$ 几何光学区，稳相近似精确，WGS 只需小幅修正
  \item $\beta < 10 \rightarrow$ 衍射区，初始相位偏差大，WGS 需要大幅修正
\end{itemize}

对于 beam=3.8mm, $f=200\text{mm}$: $\beta_x=14.8$, $\beta_y=5.4$。X 方向在几何光学区，Y 方向在衍射区边缘。

二维相位是可分离的: $\phi(x,y) = \phi_x(x) + \phi_y(y)$。

% ======================================================================
\section{DOE 面的输入光场}

\textbf{文件}: \texttt{src/propagation.py}

\begin{lstlisting}
# src/propagation.py, line 48-66
def make_input_gaussian(shape, dx_doe_m, gaussian_1e2_diameter_m,
                        clear_aperture_m, xp, dtype):
    """构建 DOE 平面高斯输入振幅."""
    Ny, Nx = shape
    # 坐标网格 (axis 0 = y, axis 1 = x)
    x = (xp.arange(Nx, dtype=dtype) - Nx//2) * dtype(dx_doe_m)
    y = (xp.arange(Ny, dtype=dtype) - Ny//2) * dtype(dx_doe_m)
    X, Y = xp.meshgrid(x, y)
    r2 = X*X + Y*Y

    w = dtype(gaussian_1e2_diameter_m / 2.0)  # 1/e^2 强度半径
    amp = xp.exp(-r2 / (w * w))               # 高斯振幅 A=exp(-r^2/w^2)

    # 圆形通光孔径截断 (直径 15mm)
    aperture_radius = dtype(clear_aperture_m / 2.0)
    amp = xp.where(r2 <= aperture_radius**2, amp, 0.0)

    # 总功率归一化到 1
    return normalize_power(amp, xp, target_power=1.0)
\end{lstlisting}

\textbf{关键点}: $A(r) = \exp(-r^2/w^2)$——注意是强度 $I$ 的 $1/e^2$ 半径，对应振幅的 $1/e$ 半径。$w = 1/e^2\text{强度直径} / 2 = 1.9\text{mm}$ (对 beam=3.8mm)。

% ======================================================================
\section{傅里叶透镜传播}

\textbf{文件}: \texttt{src/propagation.py}

DOE 面和焦面之间通过一个傅里叶透镜耦合。前向传播 = 从 DOE 到焦面（我们看到的），反传 = 从焦面回到 DOE（迭代中用来更新相位）。

\begin{lstlisting}
# src/propagation.py, line 13-16
def forward_fft(field_in, xp):
    """DOE 平面 -> 焦平面: FFT."""
    # ifftshift 把零频移到 (0,0) -> fft2 -> fftshift 把零频移回中心
    return xp.fft.fftshift(xp.fft.fft2(xp.fft.ifftshift(field_in), norm="ortho"))

def backward_fft(field_out, xp):
    """焦平面 -> DOE 平面: 逆 FFT."""
    return xp.fft.fftshift(xp.fft.ifft2(xp.fft.ifftshift(field_out), norm="ortho"))
\end{lstlisting}

\textbf{关键约定}:
\begin{itemize}[nosep]
  \item \texttt{norm="ortho"}: 保证 FFT 前后总功率守恒
  \item 矩阵方向: \texttt{axis 0 = y (行)}, \texttt{axis 1 = x (列)}——和 MATLAB \texttt{meshgrid} 一致
  \item 焦面坐标: $x_{\text{focal}}[k] = (k - N/2) \times \text{focal\_dx}$, 其中 $\text{focal\_dx} = 2.5\mu\text{m}$
\end{itemize}

\textbf{焦面像素间距}: $\text{focal\_dx} = \lambda \times f / (N \times \text{dx\_doe}) = 2.5\mu\text{m}$, $2048\times2048$ 的焦面覆盖 $\pm 2.56\text{mm}$。

% ======================================================================
\section{MRAF 焦面投影 — 三个区域，三种规则}

\textbf{文件}: \texttt{src/mraf\_gs.py}

这是每次迭代里最关键的一步: 把计算出的焦面复振幅"投影"到物理约束上。

\begin{lstlisting}
# src/mraf_gs.py, line 287-326 (精简, 去掉已注释的 GS 分支)
def _project_farfield(farfield, target_amp_eff, masks_b, xp, method,
                       mraf_factor, bg_mode, bg_factor):
    """MRAF 焦面投影."""
    # 提取当前焦面相位
    phase_ff    = xp.angle(farfield)
    phase_factor = xp.exp(1j * phase_ff)
    signal = masks_b["mask_signal"]

    projected = farfield.copy()
    # (1) signal 区: 保留相位, 替换振幅为目标振幅
    projected[signal] = target_amp_eff[signal] * phase_factor[signal]

    # (2) free 区: 保留相位, 振幅 x mraf_factor (0.4)
    free = masks_b["mask_free"]
    projected[free] *= mraf_factor

    # (3) 远背景: keep / attenuate / zero
    bg = masks_b.get("mask_bg_far", ...)
    if bg_mode == "attenuate":
        projected[bg] *= bg_factor        # x 0.9, 几乎等同 keep
    elif bg_mode == "zero":
        projected[bg] = 0

    return projected
\end{lstlisting}

\textbf{三区处理总结}:

\begin{center}
\begin{tabular}{p{0.18\textwidth}p{0.25\textwidth}p{0.18\textwidth}p{0.28\textwidth}}
\toprule
\textbf{区域} & \textbf{振幅处理} & \textbf{相位处理} & \textbf{效果} \\
\midrule
\texttt{mask\_signal} & 替换为 target\_amp $\times$ WGS\_weight & 保留原相位 & 强制趋近目标平顶 \\
\texttt{mask\_free}   & $\times 0.4$ (衰减) & 保留原相位 & 给噪声自由度，不完全压制 \\
\texttt{mask\_bg\_far} & $\times 0.9$ (几乎不变) & 保留原相位 & 不做强约束，避免能量倒灌 \\
\bottomrule
\end{tabular}
\end{center}

\textbf{为什么 \texttt{mraf\_factor}=0.4 而不是 0}: 如果 free 区完全置零，能量无处可去，会回流到 signal 区形成波纹。给 free 区一定自由度反而让平顶更平。

% ======================================================================
\section{WGS 权重更新 — 平顶内的局部修匀}

\textbf{文件}: \texttt{src/mraf\_gs.py}

MRAF 只做区域级约束。WGS 在 \texttt{mask\_flat} 内部做像素级权重反馈，进一步压平不均匀。

\begin{lstlisting}
# src/mraf_gs.py, line 164-213 (精简)
def _update_flat_wgs_weights(weights, farfield_amp, update_region,
                              xp, backend, feedback_exponent,
                              clip_min, clip_max, normalize_weights):
    """WGS flat_local 权重更新."""
    eps = np.float32(1e-12)

    # 1. 提取 mask_flat 内的焦面振幅
    amp_flat = farfield_amp[update_region]

    # 2. 计算 flat 区平均振幅
    amp_mean = xp.mean(amp_flat)

    # 3. 核心公式: 每个像素的修正系数
    #    ratio = mean(|E_flat|) / |E_pixel|
    #    如果某像素比平均暗 -> ratio > 1 -> 权重增大 -> 下次投影时要求更高
    #    如果某像素比平均亮 -> ratio < 1 -> 权重减小 -> 下次投影时要求更低
    ratio = amp_mean / xp.maximum(amp_flat, eps)
    updated = weights[update_region] * xp.power(ratio, feedback_exponent)

    # 4. Clip: 防止权重跑飞 (默认 0.5 ~ 1.5)
    updated = xp.clip(updated, clip_min, clip_max)

    # 5. 均值归一化: 保持所有权重的平均值为 1.0
    if normalize_weights:
        updated = updated / xp.mean(updated)

    weights[update_region] = updated
    return weights
\end{lstlisting}

\subsection{直观理解}

假设 mask\_flat 内某个像素振幅是 mean 的 0.8 倍，exponent=0.8:

\[
\text{ratio} = 1.0 / 0.8 = 1.25
\]
\[
\text{weight\_new} = \text{weight\_old} \times 1.25^{0.8} = \text{weight\_old} \times 1.20
\]

这个像素的权重增加了 20\%。下次 MRAF 投影时，该像素被要求匹配更高目标振幅 $\rightarrow$ 迭代推动它变亮 $\rightarrow$ 趋近平坦。

% ======================================================================
\section{主迭代循环 — 完整流程}

\textbf{文件}: \texttt{src/mraf\_gs.py}, \texttt{run\_refinement()} 函数

每轮迭代做 5 件事:

\begin{lstlisting}
# src/mraf_gs.py, line 538-609 (精简, 去掉注释和日志)
for it in range(1, total_iters + 1):          # total_iters = 200

    # (1) 前向传播: DOE -> 焦面
    field   = input_amp * exp(1j * phase)     # DOE 面复振幅
    farfield = forward_fft(field)             # FFT -> 焦面复振幅
    farfield_amp = abs(farfield)              # 焦面振幅

    # (2) WGS 权重更新 (每 5 轮一次)
    if it % update_every == 0:
        weights = _update_flat_wgs_weights(   # 只在 mask_flat 内更新
            weights, farfield_amp, update_region,
            exponent=0.8, clip=[0.5, 1.5], normalize=True)

    # (3) 加权目标振幅
    target_weighted = base_target.copy()
    target_weighted[mask_flat] *= weights     # flat 区 x WGS 权重

    # (4) MRAF 投影 + 反传 + 相位提取
    projected  = _project_farfield(           # 三区投影
        farfield, target_weighted,
        mraf_factor=0.4, bg_factor=0.9)
    nearfield  = backward_fft(projected)      # IFFT -> DOE 面
    phase      = angle(nearfield) % 2pi       # 提取新相位

    # (5) 指标记录 (每 10 轮)
    if it % metrics_interval == 0:
        记录 RMS, size50, 效率...
\end{lstlisting}

\textbf{数据流示意图}:

\begin{center}
\begin{tabular}{ccccc}
\multicolumn{2}{c}{\textbf{DOE 面}} & & \multicolumn{2}{c}{\textbf{焦面}} \\
\cline{1-2} \cline{4-5}
$\phi_{\text{old}}$ & & (1) FFT & & \\
$A_{\text{input}}$ & $\xrightarrow{\hspace{2.5em}}$ & $\longrightarrow$ & $\xrightarrow{\hspace{2.5em}}$ & $E_{\text{farfield}}$ \\
& & & & $\downarrow$ (2) WGS权重更新 \\
& & & & $\downarrow$ (3) 加权目标 \\
& & (4) IFFT & & $\downarrow$ (4) MRAF投影 \\
$\phi_{\text{new}}$ & $\xleftarrow{\hspace{2.5em}}$ & $\longleftarrow$ & $\xleftarrow{\hspace{2.5em}}$ & \\
\cline{1-2} \cline{4-5}
\end{tabular}
\end{center}

% ======================================================================
\section{ROM $\rightarrow$ SLM — 复振幅插值}

\textbf{文件}: \texttt{convert\_to\_slm.py}

计算用的 $2048\times2048$ 相位不能直接送到 SLM ($1024\times1024$)。需要裁切 + 缩放到 SLM 分辨率和像素间距。

\begin{lstlisting}
# convert_to_slm.py, line 84-106 (精简)
# 1. 根据 SLM 物理尺寸 (17.4mm) / 计算网格像素间距 (20.8um) -> 裁切 837x837 像素
phys_mm  = 1024 * 17.0e-3                     # SLM 物理尺寸 = 17.4mm
crop_px  = int(phys_mm / (dx_doe_um * 1e-3))  # 裁切像素数 = 837
phase_crop = phase[cy-418:cy+419, cx-418:cx+419]  # 从 2048 中心裁出

# 2. * 关键: 复振幅插值 (不是直接插值相位!)
zoom_ratio = 1024 / 837                        # ~ 1.22x 放大
cfield = np.exp(1j * phase_crop)               # 转为复振幅
real_z = zoom(cfield.real, zoom_ratio, order=3)  # 实部 cubic 插值
imag_z = zoom(cfield.imag, zoom_ratio, order=3)  # 虚部 cubic 插值
phase_slm = np.arctan2(imag_z, real_z)          # 从插值后的复振幅恢复相位
phase_slm = phase_slm % (2*pi)
\end{lstlisting}

\textbf{为什么不能直接插值包裹相位}: 包裹相位在 $0 \leftrightarrow 2\pi$ 处有跳变，cubic spline 会把跳变当成高频信号产生巨大振铃。先转成连续平滑的 $\exp(i\phi)$ 分量，插值完再恢复相位，完全避免这个问题。

% ======================================================================
\section{闪耀光栅 — 让光斑偏离中心}

在 DOE 相位上叠加线性斜坡，焦面光斑就会平移。

\begin{lstlisting}
# 光栅相位: phi_grating(x) = 2*pi * (dx / focal_dx) * x / N
focal_dx_um = 2.5      # 焦面像素间距
shift_x_um  = 400.0    # 想偏移 400um
shift_y_um  = 200.0    # 想偏移 200um

periods_x = shift_x_um / focal_dx_um   # = 160 个周期 (跨 2048 像素)
periods_y = shift_y_um / focal_dx_um   # = 80  个周期

x = np.arange(N) - N//2
X, Y = np.meshgrid(x, x)
grating = 2*pi * (periods_x * X + periods_y * Y) / N

phase_with_grating = (phase_doe + grating) % (2*pi)
\end{lstlisting}

\textbf{物理原理}: 线性相位斜坡 = 平面波倾斜入射 = 焦面光斑侧移。1 个光栅周期（跨整个 DOE 孔径）= 焦面偏移 1 个像素 = $2.5\mu\text{m}$。

% ======================================================================
\section{指标计算 — 怎么判断做得好不好}

\textbf{文件}: \texttt{src/metrics.py}

\begin{lstlisting}
# src/metrics.py, line 86-122 (精简)
def compute_metrics(intensity, masks, x_um, y_um, iteration):
    I = intensity
    mask_flat = masks["mask_flat"]

    # (1) 归一化: I_norm = I / mean(I[mask_flat])
    #    不以全局最大值为基准, 因为那会被偶然的 hot pixel 拉偏
    I_n = I / np.mean(I[mask_flat])

    # (2) RMS 不均匀度 (最核心指标)
    flat_vals = I_n[mask_flat]
    flat_rms  = np.std(flat_vals) / np.mean(flat_vals)
    # 例如: std=0.0087, mean=1.0 -> RMS = 0.87%

    # (3) size50: 中心剖面半高全宽
    profile_x = I_n[N//2, :]
    size50_x  = width_at_level(x_um, profile_x, level=0.5)
    # 从中心向两侧找到 I=0.5 的 crossing 点, 线性插值

    # (4) 效率: signal 区内能量 / 总能量
    efficiency = np.sum(I[mask_signal]) / np.sum(I)

    return { "flat_rms": flat_rms, "size50_x": size50_x, ... }
\end{lstlisting}

\textbf{RMS 定义}: $\operatorname{std}(\text{mask\_flat}) / \operatorname{mean}(\text{mask\_flat})$，只算 flat core 内的像素。对于 beam=3.8mm，200 轮后 RMS ≈ 0.71\%。

% ======================================================================
\section{参数速查表}

\begin{center}
\begin{tabular}{p{0.35\textwidth}p{0.22\textwidth}p{0.33\textwidth}}
\toprule
\textbf{参数} & \textbf{值} & \textbf{含义} \\
\midrule
$\lambda$                              & 532 nm          & 激光波长 \\
$f$                                    & 200 mm          & 透镜焦距 \\
$N$                                    & 2048            & 计算网格 \\
focal\_dx                              & 2.5 $\mu$m      & 焦面像素间距 \\
beam (实测)                            & 3.8 mm          & $1/e^2$ 强度直径 \\
目标 $W_{50}\times H_{50}$            & $330\times120$ $\mu$m & 平顶 50\% 尺寸 \\
$\delta_x$ / $\delta_y$                & 15 / 8 $\mu$m   & 边缘过渡半宽 \\
mraf\_factor                           & 0.4             & free 区衰减系数 \\
bg\_factor                             & 0.9             & 背景衰减 (接近 keep) \\
feedback\_exponent                     & 0.8             & WGS 修正力度 \\
weight\_clip                           & [0.5, 1.5]      & 权重裁剪范围 \\
update\_every                          & 5               & 每 N 轮更新一次权重 \\
num\_iters                             & 200             & 总迭代轮数 \\
\bottomrule
\end{tabular}
\end{center}

% ======================================================================
\section{文件索引}

\begin{center}
\begin{tabular}{p{0.38\textwidth}p{0.52\textwidth}}
\toprule
\textbf{想了解什么} & \textbf{去看哪个文件} \\
\midrule
初始相位怎么来的         & \texttt{make\_phase0.py} (RD 公式) \\
目标光斑怎么定义         & \texttt{src/rtad\_target.py} (升余弦 + mask) \\
FFT 传播怎么写的         & \texttt{src/propagation.py} \\
迭代循环逻辑             & \texttt{src/mraf\_gs.py} $\rightarrow$ \texttt{run\_refinement()} \\
焦面投影公式             & \texttt{src/mraf\_gs.py} $\rightarrow$ \texttt{\_project\_farfield()} \\
WGS 权重怎么更新         & \texttt{src/mraf\_gs.py} $\rightarrow$ \texttt{\_update\_flat\_wgs\_weights()} \\
怎么转成 SLM 格式        & \texttt{convert\_to\_slm.py} \\
光斑质量怎么量化         & \texttt{src/metrics.py} $\rightarrow$ \texttt{compute\_metrics()} \\
CLI 参数怎么传           & \texttt{run\_rtad\_mraf\_gs\_case.py} $\rightarrow$ \texttt{parse\_args()} \\
GPU/CPU 怎么选           & \texttt{src/backend.py} \\
MATLAB 文件怎么读        & \texttt{src/io\_mat.py} \\
\bottomrule
\end{tabular}
\end{center}

\end{document}
